%! Question groups, for make test (tests/check_groups.py).
% Built once as written, once as a key, and once for each broken case, chosen
% by defining \GroupsCase in the pretex: undeclared, empty, duplicate, points.
% Each broken case must stop the build with its error.
\documentclass[10pt]{article}
\usepackage{example-article}
\gradingassignment{test-groups}
\providecommand\GroupsCase{ok}
\makeatletter
\def\GroupsIf#1#2{\edef\@tempa{#1}\edef\@tempb{\GroupsCase}%
  \ifx\@tempa\@tempb #2\fi}
\makeatother

\begin{document}

A two-part question with the default rubric path, its parts written before
and after the declaration. \gradedblank[id=q1a, type=rubric, question=q1]{\linewidth}{1in}
\gradedquestion[id=q1, scheme=ap-stats]
\gradedblank[id=q1b, type=short, question=q1, points=3, answer=7]{1in}{0.3in}

A question with its own rubric path and no scheme written.
\gradedquestion[id=q2, rubric={../rubrics/test-groups-q2.yaml}]
\gradedblank[id=q2a, type=choice, question=q2, answer=B]{0.4in}{0.4in}
\gradedblank[id=q2b, type=rubric, question=q2, points=0]{\linewidth}{1in}

An ordinary item: \gradedblank[id=e1, type=short, points=2, answer=5]{1in}{0.3in}

\GroupsIf{undeclared}{\gradedblank[id=z1, type=short, question=zz, answer=1]{1in}{0.3in}}
\GroupsIf{empty}{\gradedquestion[id=q9]}
\GroupsIf{duplicate}{\gradedquestion[id=q1, rubric=other.yaml]}
\GroupsIf{points}{\gradedquestion[id=q8, points=5]}

\end{document}
